\scanpage{52}
\mspage{1}

Rétrospection sur les cartes. (finies)

\textcircled{1} Définition initiale \inserted{d'une \uline{carte} (topologique) cellulaire}: surface compacte (pas néc. orientable) $X$, avec partie fermée $K$, loc. \unsure{homéom.} à une étoile (``1-complexe topologique''), plus hypothèses ``cellulaires'': a) $K$ n'a pas de composantes connexes homéom. à $S^1$ \deleted{\unsure{ni de}} \inserted{ni de points isolés} b) Les comp. connexes de $X-K$ sont des disques ouverts i.e. homéom. à $\mathring D=\{z\in\mathbb C\mid |z|<1\}$.

\margin{b) implique déjà que $K$ n'a pas de pt isolé, \uline{sauf} \deleted{\unsure{quand}} il y a une comp. connexe $X_0$ de $X$ \deleted{\unsure{réduite qui}} est une sphère, avec $K\cap X_0$ réduit à un point.}

On en fait une ``catégorie isotopique'', et à l'aide de la notion de \unsure{biarc} (p.ex.), on établit une équivalence de catégories
\[
(1)\qquad \deleted{\text{\unsure{Cartes}}}\quad\text{Cartes isot.}\ \underset{\rule{1.1cm}{0.3pt}}{\approx}\ G_2\text{-ensembles }\inserted{\text{finis}}\text{ ``spéciaux''}
\]
où
\[
(2)\qquad G_2\simeq\{\tau_0,\tau_1,\underbrace{\tau_2}_{\text{ou }\tau_\infty}\mid
\tau_0^2=\tau_1^2=\tau_2^2=(\tau_0\tau_2)^2=1\}
\]
(\uline{NB} Par la suite, on notera $\unsure{\sigma_\infty|\tau_\infty}$ en lieu de $\unsure{\sigma_2|\tau_2}$, notation mieux appropriée pour les relations avec les Cour. Alg. La notation $\unsure{\sigma_2|\tau_2}$ est plus appropriée pour les généralisations en dim supérieure $\ldots$),
\margin{La notation $\unsure{\sigma_2|\tau_2}$ est plus commode p.ce. qu'on dira que $\unsure{\sigma_i|\tau_i}$ est l'opération (sur les repères) qui \unsure{change le $i$-feuillet du repère en le} \unsure{de façon à fixer les autres}.}
et où les $G$-ens ``spéciaux'' sont ceux où les involutions
\[
\tau_0,\tau_1,\tau_2\quad\text{et}\quad\sigma=\tau_0\tau_2
\]
opèrent sans points fixes. (NB. Par la suite, on pose $\rho_1=\tau_0\tau_\infty$, en $=$ temps qu'on écrit $\tau_\infty$ pour $\tau_2$.)

On introduit
\[
(3)\qquad G_2^\circ=\operatorname{Ker}\left(\begin{array}{rcl}
G_2&\longrightarrow&\mathbb Z/2\mathbb Z\\
\tau_i&\longmapsto&1\bmod2
\end{array}\right)
\]
et on a
\[
\unsure{(4)}\qquad G_2^\circ=\{\rho_0,\rho_1,\rho_\infty\mid\rho_\infty\rho_1\rho_0=\deleted{\unsure{1}}\rho_1^2=1\}
\]
\mspage{51}
